\clearpage
\VerdierPhysicalPageAnchor{0238}%
\begingroup\small\setlength{\parskip}{0pt}\setlength{\emergencystretch}{1.5em}
\setlength{\abovedisplayskip}{4pt plus 2pt minus 2pt}%
\setlength{\belowdisplayskip}{4pt plus 2pt minus 2pt}%

\VerdierSemanticAnchor{ch4d-4-3-4}
\noindent\textbf{4.3.4.}\quad
The functor from $D(\mathcal A)$ to the category of spectral objects of
$D(\mathcal A)$ is still called, by an abuse of terminology, the
\emph{second canonical spectral-object functor}.

Let $X$ be a complex in $\mathcal A$, and let
$\bigl(X(p,q),\delta(p,q,r)\bigr)$ be the second spectral object associated
with $X$. It follows from (4.3.2) and (1.2.10) that there is a canonical
isomorphism in $D(\mathcal A)$:

\VerdierSemanticAnchor{ch4d-4-3-4-1}
\begin{equation}
X(p,p+1)\xrightarrow{\sim}H^p(X)[-p].
\tag{4.3.4.1}
\end{equation}

\VerdierSemanticAnchor{ch4d-4-3-4-boundedness}
Notice also that when $X$ is an object of $D^*(\mathcal A)$
($*=-,+,b,$ or ``blank''), the second spectral object associated with $X$
in fact has values in $D^*(\mathcal A)$.

\VerdierSemanticAnchor{ch4d-4-4}
\medskip\noindent\textbf{4.4. Usual spectral sequences.}

\VerdierSemanticAnchor{ch4d-4-4-1}
\medskip\noindent\textbf{4.4.1.}\quad
Let $\mathcal A$ be an additive category, not necessarily abelian, let
$F\colon K(\mathcal A)\longrightarrow\mathcal B$ be a cohomological
functor into an abelian category $\mathcal B$, and let
$F^p\colon K(\mathcal A)\longrightarrow\mathcal B$ be the functors
obtained from $F$ by composition with the translation functor of
$K(\mathcal A)$ ($F^p(X)=F(X[p])$). In addition, let $Y$ be a complex in
$\mathcal A$. The functor $F$ takes the first spectral object with values
in $K(\mathcal A)$ associated with $Y$ (4.1.4) to a spectral object with
values in $\mathcal B$. This yields (Chapter~II, 4.3.3) a spectral sequence
with values in $\mathcal B$, functorial in $F$ and in $Y$ (when the latter
varies in $\operatorname{comp}(\mathcal A)$), called the \emph{first
spectral sequence of the functor $F$ relative to the complex $Y$}. For
this spectral sequence we always use the indexing introduced in
(Chapter~II, 4.3.3.1).

\VerdierSemanticAnchor{ch4d-4-4-2}
\medskip\noindent\textbf{Proposition 4.4.2.}\quad

\VerdierSemanticAnchor{ch4d-4-4-2-a}
\noindent\textit{a) The term $\mathrm I_1^{p,q}$ of the first spectral
sequence of the functor $F$ relative to the complex $Y$ is given by}
\[
\mathrm I_1^{p,q}=F^q(Y^p)
\]
\noindent\textit{(here $Y^p$ denotes the complex whose only nonzero
component is, in degree zero, the $p$-th component of the complex $Y$).
The corresponding differential is}

\VerdierSemanticAnchor{ch4d-4-4-2-a-differential}
\[
\begin{aligned}
d_1&\colon\mathrm I_1^{p,q}\longrightarrow\mathrm I_1^{p+1,q},\\
d_1&=F^q(-d_Y^p)\colon F^q(Y^p)\longrightarrow F^q(Y^{p+1}).
\end{aligned}
\]

\VerdierSemanticAnchor{ch4d-4-4-2-b}
\noindent\textit{b) The term $\mathrm I_2^{p,q}$ of this same spectral
sequence can therefore be described as follows: denote by
$F^{q,\bullet}\colon\operatorname{comp}(K(\mathcal A))
\longrightarrow\operatorname{comp}(\mathcal B)$ the extension of $F^q$
to complexes, by $H_{\mathcal B}^p\colon\operatorname{comp}(\mathcal B)
\longrightarrow\mathcal B$ the $p$-th cohomology-object functor, and by
$Y^\bullet$ the following object of $\operatorname{comp}(K(\mathcal A))$:}

\VerdierSemanticAnchor{ch4d-4-4-2-b-complex}
\[
Y^\bullet=\cdots\longrightarrow Y^{p-1}
\xrightarrow{d_Y^{p-1}}Y^p\xrightarrow{d_Y^p}Y^{p+1}
\longrightarrow\cdots.
\]
\endgroup
